\section{Introduction}\label{sec:intro}

A common way to commit to a multivariate polynomial
$f\in\F[x_1,\dots,x_m]$ whose individual degrees are below $k$, written
$f\in\F[x_1,\dots,x_m]_{<k}$, is to commit to one
univariate polynomial with the same coefficients, its Kronecker image
\[
U_f(X)=f\big(X,X^{k},\dots,X^{k^{m-1}}\big),
\]
and to run a Reed--Solomon proximity test such as FRI~\cite{BBHR18} on its evaluations. The
coefficient vector of $f$ is the coefficient vector of $U_f$, the commitment is a
Reed--Solomon codeword, and the deployed machinery for univariate polynomials applies
unchanged. Multilinear commitment schemes built on FRI follow this route~\cite{ZCF24}, and so
do schemes for polynomials of higher individual degree~\cite{PolyFRIM}.

What such a test certifies is a statement about $U_f$: its degree is below $k^m$. Through
the Kronecker map this means that every individual degree of $f$ is below $k$, and nothing
more. The total degree of $f$, or a weighted combination of its degrees, does not appear in
the degree of $U_f$; the polynomials $x_1$ and $x_1^{k-1}\cdots x_m^{k-1}$ both pass the degree test on $U_f$. Yet
total-degree bounds are the classical object of multivariate low-degree
testing~\cite{RS96}, and proximity tests for them exist only for codes evaluated on product
sets $L^m$~\cite{ABN23,MZ25}, not for the univariate encoding that FRI-based systems use.

We give such a test. It is FRI with one change: the folding challenges are symbolic. In
round $i$ the verifier sends $r_i\in\F$ and the fold uses $r_iT^{w_i}$, where $T$ is a formal
variable and $w=(w_1,\dots,w_m)\in\N^m$ a fixed weight vector. Write
$f=\sum_{\alpha}f_\alpha x^\alpha$, where $\alpha=(\alpha_1,\dots,\alpha_m)$ ranges over
$\{0,\dots,k-1\}^m$, $x^\alpha:=x_1^{\alpha_1}\cdots x_m^{\alpha_m}$ and $f_\alpha\in\F$.
The \emph{weighted degree} of $f$ is
\[
\deg_wf:=\max\{\langle w,\alpha\rangle:\ f_\alpha\ne0\},\qquad
\langle w,\alpha\rangle:=w_1\alpha_1+\dots+w_m\alpha_m,
\]
with $\deg_w0:=-\infty$. Since a fold of a Kronecker image substitutes the
challenge for one variable, the honest last message is
\[
g(T)=f\big(r_1T^{w_1},\dots,r_mT^{w_m}\big),
\]
whose degree in $T$ equals $\deg_wf$ for all but a fraction at most $m(k-1)/|\F|$ of the
challenges $(r_1,\dots,r_m)\in\F^m$. The verifier
checks $\deg g\le D$.

\paragraph{Why the univariate encoding.} A FRI-based system that commits to $f$ through
$U_f$ already holds the word on which our test runs. Testing a degree bound on $f$ therefore
needs neither a second commitment nor an evaluation domain of the form $L^m$, and the
verifier keeps the structure of FRI: one fibre per level, one consistency check per fibre.
The weighted degree covers the total degree
($w=(1,\dots,1)$, so $\langle w,\alpha\rangle=\alpha_1+\dots+\alpha_m$) and the degree in a
single variable $x_l$ ($w=e_l$, the $l$-th unit vector of $\N^m$, so
$\langle w,\alpha\rangle=\alpha_l$) with one definition and one analysis.

\paragraph{Why symbolic challenges.} The classical way to measure the total degree of $f$ is
to restrict it to a random line and read off the degree of the univariate
restriction~\cite{RS96}. On a Kronecker commitment this would require evaluations of $f$ at
arbitrary points of $\F^m$, that is, opening proofs on top of the proximity test. The symbolic
fold produces such a restriction without any opening: $g(T)$ is the restriction of $f$ to the
curve $T\mapsto(r_1T^{w_1},\dots,r_mT^{w_m})$, and it is the last message of the protocol.

\paragraph{What it costs.} The oracles are $D+1$ times larger than those of FRI at most, and a
query reads at most $D+1$ times more field elements: a value at level $i$ is a polynomial in
$T$ of degree at most $d_i=\min\big(D,(k-1)\sum_{l\le i}w_l\big)$. The number of rounds, the
number of query repetitions for a given soundness, and the shape of the verifier are those of
FRI.

\subsection{Results}\label{sec:intro-results}

For $w\in\N^m$ and $D\in\N$, the \emph{degree code} $C_{w,D}$ is the set of evaluations on
$L_0$ of the Kronecker images $U_f$ with $f\in\F[x_1,\dots,x_m]_{<k}$ and $\deg_wf\le D$
(Definition~\ref{def:graded}). Our main result is an interactive oracle proof of proximity
for $C_{w,D}$ that reaches the Johnson radius.

\begin{introthm}[informal]\label{thm:intro}
Let $\F$ be a finite field, $k\ge2$, $m\ge1$, $N$ an integer with $k^m\mid N\mid|\F|-1$,
$L_0\subseteq\F^*$ a coset of the subgroup of order $N$ of $\F^*$, and $\rho:=k^m/N$. For every $w\in\N^m$ and $D\in\N$ there is a
public-coin IOPP $\Pi_{w,D}$ for $C_{w,D}$ with the following properties.
\begin{itemize}
\item \emph{Rounds and messages:} $m$ rounds; the message of round $i$ is a word on a domain
of size $N/k^i$ whose values are polynomials in $T$ of degree at most
$d_i=\min\big(D,(k-1)\sum_{l\le i}w_l\big)$.
\item \emph{Completeness:} perfect.
\item \emph{Soundness:} if the committed word has agreement at most $\alpha$ with
$C_{w,D}$ (the largest fraction of $L_0$ on which it coincides with a codeword,
Definition~\ref{def:agree}), where $\alpha>\sqrt\rho$, the verifier accepts with probability at most
$\varepsilon+\alpha^s$, with $\varepsilon=O\big(N^2/|\F|\big)$ for fixed $\rho$ and
$\alpha-\sqrt\rho$, and $s$ the number of query repetitions.
\item \emph{Query cost:} each repetition reads $k\sum_{i=1}^m(d_{i-1}+1)\le km(D+1)$ field
elements.
\end{itemize}
\end{introthm}

The formal statement is \cref{thm:main}, a reduction to the hypothesis of \cref{sec:params};
\cref{cor:johnson} gives the explicit constants at the Johnson radius and \cref{cor:ud} a simpler
error in the unique-decoding regime. With $w=(1,\dots,1)$ the test certifies the
total degree; with $w=e_l$ and $D=b$ it certifies $\deg_{x_l}f\le b$ for any $b<k$, which the
Kronecker rate alone cannot do (\cref{sec:individual}). Several constraints can be tested in
one run with several formal variables (\cref{rem:multi}). At $N=2^{20}$, $\rho=2^{-4}$,
$k=4$, $m=8$, $D=6$ and $|\F|=2^{256}$, $65$ repetitions give soundness error below
$2^{-129.9}$ against words with agreement at most $\alpha_0:=\sqrt\rho\,(1+\frac1{2\kappa})\approx0.25006$
with $C_{w,D}$, where $\kappa=2^{11}-1$ is the parameter of \cref{cor:johnson},
that is, at distance just inside the Johnson radius. This is the number of repetitions FRI
needs at the same distance~\cite[\S8.2]{BCIKS20}, and each repetition reads $188$ field
elements against $32$ for FRI (\cref{sec:numbers}). After the commitment the prover performs
$O(k^2(D+1)N/(k-1))$ field operations with direct interpolation on each fibre.

\subsection{Correlated agreement for subcodes}\label{sec:intro-ca}

The analysis follows the proof of FRI soundness by Ben-Sasson, Carmon, Ishai, Kopparty and
Saraf~\cite{BCIKS20}, which rests on correlated agreement over curves: if many points of a
curve $\sum_jz^ju_j$ are close to a Reed--Solomon code, then the words $u_j$ agree with
codewords on a common large set. Two features of our setting are not covered by that
statement. The words take values in $\F[T]$, and the target codes are subcodes of
Reed--Solomon codes. The first is handled by reading $\F[T]$-valued words as words over an
extension field $\E$ while the challenges stay in $\F$; this relies on a property of the
proofs in~\cite{BCIKS20}, that they use the set of good challenges only through its size and
membership, which we record with page references (\cref{sec:params}). The second is handled
by a short linear-algebra lemma.

\begin{introthm}[informal; \cref{cor:share}]
Weighted correlated agreement over curves whose parameter lies in $\F$ holds, with the same
parameters, for every $\F$-linear subcode of a Reed--Solomon code over an extension of $\F$:
the words $u_j$ agree on a common large set with codewords of the subcode.
\end{introthm}

The statement is not specific to degrees; it applies, for instance, to the constrained
Reed--Solomon codes of WHIR~\cite{ACFY25}. It also makes our soundness theorem a reduction:
any weighted correlated-agreement theorem for Reed--Solomon codes that discharges the
hypothesis \cref{hyp:ca} transfers to the degree test at the same radius. The weighted form is
needed because the analysis of FRI that we follow~\cite[\S8]{BCIKS20} is weighted. Beyond the
Johnson radius, correlated agreement over curves up to radius $1-\rho-\gamma$ has recently
been claimed for prime fields~\cite{Jer26}, in an unweighted form. A weighted version, stated
as follows, would extend \cref{thm:main} accordingly.

\begin{conjecture}\label{conj:capacity}
For every $\gamma>0$ and $l\ge1$ there is $c=c(\gamma,l)$ such that
$\mathrm H(\rho+\gamma,\Theta)$ holds for $\Theta(n,M,l):=l\,(Mn+1)^c$.
\end{conjecture}

For $\alpha\ge\alpha_0$ the hypothesis is \cref{prop:H1}, which is what \cref{cor:johnson}
uses. If \cref{conj:capacity} holds, \cref{thm:main} applies verbatim and $\Pi_{w,D}$ is sound
for agreement $\alpha\ge\rho+\gamma$, with error at most $m(k-1)(N+1)^{c}/|\F|+\alpha^s$.

\subsection{Comparison}\label{sec:intro-comparison}

\Cref{tab:comparison} places the test among proximity tests for multivariate degree.

\begin{table}[t]
\centering\footnotesize
\renewcommand{\arraystretch}{1.2}
\begin{tabular}{@{}p{3.1cm}p{2.1cm}p{2.6cm}p{2.9cm}p{3.0cm}@{}}
\hline
 & oracle & degree certified & radius & reads per repetition\\
\hline
FRI on $U_f$ \cite{BBHR18,BCIKS20} & $U_f$ on $L_0$ & individual, $<k$ & Johnson & $km$ field elements\\
PolyFRIM \cite{PolyFRIM} & coefficient vector & individual, $<k$ & unique decoding & as FRI\\
tensor-RS IOPP \cite{ABN23} & $f$ on $L^m$ & individual & $1-\sqrt[3]{1-\lambda}$ & about $2^m\log|L|$ points\\
Reed--Muller IOPP \cite{ABN23} & $f$ on $L^m$ & total & $1-\sqrt[3]{1-\lambda}$ & about $2^m\log|L|$ points\\
\textbf{this work} & $U_f$ on $L_0$ & weighted, total, individual & Johnson & $\le km(D+1)$ field elements\\
\hline
\end{tabular}
\caption{Proximity tests for multivariate degree. ``Radius'' is the proximity range covered
by the stated soundness analysis, up to slack terms; $\lambda$ is the relative minimum
distance of the code tested in~\cite{ABN23}.}
\label{tab:comparison}
\end{table}

\begin{itemize}
\item The IOPPs of~\cite{ABN23} test individual and total degree for
evaluations on a product set $L^m$, the natural encoding for constructions based on
Reed--Muller codes; their total-degree folding uses two challenge vectors and balancing
terms per round, and each round reads $2^m$ points, so a repetition reads about
$2^m\log|L|$ points. We test a different code: the same
degree constraints on the univariate encoding, with $m$ rounds of $k$-ary folding.
\item \emph{PolyFRIM}~\cite{PolyFRIM} commits to multivariate polynomials through the same
coefficient vector and proves evaluations; the degree it certifies is the individual degree
fixed by the encoding. Our test adds total and weighted degree bounds on that commitment.
\item \emph{FRI} on $U_f$~\cite{BBHR18,BCIKS20} is cheaper: our test reads up to $D+1$ times
more field elements per query and sends oracles up to $D+1$ times larger. In exchange it
certifies the weighted degree, where FRI certifies only $\deg_{x_l}f<k$.
\end{itemize}

\paragraph{Organisation.} \Cref{sec:overview} explains the protocol and the proof.
\Cref{sec:prelim} fixes notation, \cref{sec:tools} proves the tools, and \cref{sec:test}
gives the protocol and its analysis. \Cref{sec:variants} treats variants,
\cref{sec:impl} describes a prototype, and \cref{sec:conclusion} lists open problems.
